\chapter*{Conclusiones}
\addcontentsline{toc}{chapter}{Conclusiones}

% Estructura propuesta. Texto continuo, sin secciones; un bloque de uno o dos
% párrafos por punto. Cada bloque dice de dónde sale su contenido.

% [1] La respuesta. La pregunta de la intro (cómo cambiaron los incidentes
% alrededor de las cámaras desde fotocívicas), el diseño en una frase (áreas
% circulares, emparejamiento y diferencias en diferencias) y el resultado: -1.8%
% a 300 m, con un intervalo de -6.0% a +2.2%, sin efecto distinguible de cero en
% ningún radio ni tipo de incidente. (Cap. 4, secciones 4.1 y 4.3.)

% [2] Contra la hipótesis. La intro plantea reducciones de 10-20%, mayores en los
% incidentes graves. Bajo tendencias paralelas, los datos descartan reducciones de
% ese tamaño (y las de 12-28% en lesionados de la literatura), pero no una de
% hasta 6% en el total o 9% en lesionados. En fatales no se puede concluir nada.
% (Cap. 4, 4.1, 4.2 y 4.3.)

% [3] Por qué difiere de lo reportado para la CDMX. Retoma los tres elementos que
% el cap. 1 dice que condicionan los resultados (línea 39 de 1.Rev.tex):
%   - la escala: comparar contra toda la ciudad da -14%; con zonas comparables
%     desaparece (escalera de la 4.1 y sensibilidad de la 4.5). Quintero et al.
%     miden alcaldías.
%   - la ubicación: las cámaras se pusieron donde había más incidentes, y por eso
%     la comparación ingenua confunde zonas peligrosas con efecto.
%   - la sanción: la sanción cívica puede restar fuerza disuasiva. Es una
%     posible explicación, no algo que se pruebe aquí.
% También el efecto canguro y el tamaño del área: el cap. 2 lo plantea y nadie lo
% retoma; los resultados por radio no muestran un efecto concentrado cerca de la
% cámara.

% [4] Qué se estima y qué no. Es el efecto local de estar cerca de una cámara
% bajo el mismo régimen de sanciones, no el efecto del programa en toda la
% ciudad. Y no habla de las 6 zonas más peligrosas, que quedaron fuera.
% (conclusiones.ipynb §7, primeros dos puntos; cap. 3, unidades excluidas.)

% [5] Limitaciones que el lector debe tener en mente. La prueba de tendencias
% previas tiene poca potencia; la inferencia es aproximada (el cap. 2 promete
% aquí las limitaciones del bootstrap: no es formalmente válido para vecino más
% cercano y las 93 unidades se toman como fijas); no hay datos de tránsito ni de
% velocidad. Las limitaciones de los datos ya están en el cap. 3, así que solo se
% remite a ellas. (conclusiones.ipynb §7.)

% [6] Qué haría falta y qué implica. Datos de aforo y velocidad dentro de la
% ciudad serían la mejora más valiosa. Para la política, cautela al atribuir a
% las cámaras las caídas que se observan en toda la ciudad. Opcional: el
% emparejamiento podría usarse para elegir dónde ampliar el programa y evaluarlo.
